\chapter{Sequence Ordering }
\hypertarget{seqorder}{}\label{seqorder}\index{Sequence Ordering@{Sequence Ordering}}
The order in which sequences are listed while unsorted in Test\+Stand is directly reflected in how they show up when that sequence file is referenced elsewhere. Therefore, there\textquotesingle{}s significance in paying attention to where new sequeces are created within a sequence file.

The main sequences all are organized in such a way that sequences of similar functionality are next to each other. Any non-\/standard sequences are added at the bottom of the sequence list. Their names are also prepended with a lower case "{}z"{} so that standard and non-\/standard sequences are still grouped together when sorting by sequence name.

Additionally, this allows for divider sequences to provide further clarity. This is seen in Test\+Sequence.seq as empty sequences with names in the format of {\ttfamily ================\+[Label]\+================} to separate Tests, Templates, and Utility sequences. You should divide the test sequences up further if applicable (Eg. Traceability, Pre-\/\+Programming, Post-\/\+Programming, etc.) as there tends to be a lot of selecting specific sequences to call in Main\+Test\+Sequence of Test\+Sequence.\+seq. 